\section*{Chapitre \ref{ch:b3:complex-spectra-magnetism} --- Spectres électroniques et magnétisme des complexes}

\begin{solution}{exo:b3:complex-spectra-magnetism:1}
D~: E$_g$ + T$_{2g}$ ($2 + 3 = 5$)~; F~: A$_{2g}$ + T$_{1g}$ + T$_{2g}$ ($1 + 3 + 3 = 7$)~; G~: A$_{1g}$ + E$_g$ + T$_{1g}$ + T$_{2g}$
($1 + 2 + 3 + 3 = 9$).
\end{solution}

\begin{solution}{exo:b3:complex-spectra-magnetism:2}
$d^1$ $^2$D $\to$ $^2$T$_{2g}$~; $d^2$ $^3$F $\to$ $^3$T$_{1g}$~; $d^3$ $^4$F $\to$ $^4$A$_{2g}$~; $d^4$ $^5$D $\to$ $^5$E$_g$~; $d^5$ $^6$S $\to$ $^6$A$_{1g}$~; $d^6$ $^5$D
$\to$ $^5$T$_{2g}$~; $d^7$ $^4$F $\to$ $^4$T$_{1g}$~; $d^8$ $^3$F $\to$ $^3$A$_{2g}$~; $d^9$ $^2$D $\to$ $^2$E$_g$.
\end{solution}

\begin{solution}{exo:b3:complex-spectra-magnetism:3}
Haut spin~: \ce{Mn^{2+}} ($n = 5$) $5.92\,\mu_B$, \ce{Fe^{2+}} (4) $4.90\,\mu_B$, \ce{Co^{2+}} (3) $3.87\,\mu_B$~; bas spin~: \ce{Fe^{2+}} et
\ce{Co^{3+}} ($t_{2g}^6$)~: 0, diamagnétiques.
\end{solution}

\begin{solution}{exo:b3:complex-spectra-magnetism:4}
$\ce{[Mn(H2O)6]^{2+}}$ (interdite de spin et par la \omterm{thm:b3:electronic-spectroscopy:laporte}{règle de Laporte}, presque incolore) $<$ $\ce{[Co(H2O)6]^{2+}}$
(interdite par la \omterm{thm:b3:electronic-spectroscopy:laporte}{règle de Laporte}, vibronique) $<$ $\ce{[CoCl4]^{2-}}$ (pas de centre de symétrie) $<$ \ce{MnO4-} (transfert
de charge permis).
\end{solution}

\begin{solution}{exo:b3:complex-spectra-magnetism:5}
8500~: $^3$A$_{2g}\to{}^3$T$_{2g}$, donc $\Delta_{\mathrm o} = \qty{8500}{cm^{-1}}$~; 14100~: $^3$T$_{1g}$(F)~; 24900~: $^3$T$_{1g}$(P).
$15B = 39\,000 - 3 \times 8500$, $B = \qty{900}{cm^{-1}}$, $\beta = 900/1056 = 0.85$.
\end{solution}

\begin{solution}{exo:b3:complex-spectra-magnetism:6}
$\Delta_{\mathrm o} = \qty{17400}{cm^{-1}}$~; en résolvant $7.5B + 1.5\Delta_{\mathrm o} - \frac12\sqrt{225B^2 - 18B\Delta_{\mathrm o} + \Delta_{\mathrm o}^2} = 24\,600$~: $B = \qty{729}{cm^{-1}}$,
$\beta = 0.79$. $\nu_3 = \qty{38500}{cm^{-1}}$ (\qty{260}{nm}), dans l'ultraviolet.
\end{solution}

\begin{solution}{exo:b3:complex-spectra-magnetism:7}
Les électrons $d$ sont davantage délocalisés sur les ions oxyde des gemmes
que sur les molécules d'eau~: les liaisons Cr--O y sont plus covalentes.
\end{solution}

\begin{solution}{exo:b3:complex-spectra-magnetism:8}
Forte ($e_g$ inégalement rempli)~: $d^9$, $d^4$ haut spin, $d^7$ bas spin. Nulle~: $d^3$ ($t_{2g}^3$, terme A) et
$d^6$ bas spin ($t_{2g}^6$). Faible~: $d^6$ haut spin ($t_{2g}^4e_g^2$, terme T).
\end{solution}

\begin{solution}{exo:b3:complex-spectra-magnetism:9}
$d^7$ octaédrique haut spin~: terme fondamental $^4$T$_{1g}$, un terme T avec une \omterm{def:b3:complex-spectra-magnetism:moment}{contribution orbitale},
moment nettement supérieur à la valeur de spin seul, $3.87\,\mu_B$. $d^7$ tétraédrique ($\equiv$ $d^3$ octaédrique)~: terme
fondamental $^4$A$_2$, moment orbital bloqué, proche de la valeur de spin
seul (un peu au-dessus, par mélange avec des termes T excités).
\end{solution}

\begin{solution}{exo:b3:complex-spectra-magnetism:10}
$\mu_{\mathrm{eff}} = 2.828\sqrt{1.87} = 3.87\,\mu_B$~: trois électrons non appariés ($S = \frac32$).
\end{solution}

\begin{solution}{exo:b3:complex-spectra-magnetism:11}
$\Delta\chi_v = 3 \times 25.0/\num{400e6} = \num{1.88e-7}$~; $\chi_{\mathrm m} = \num{1.88e-7}/\qty{10.0}{mol.m^{-3}} = \qty{1.88e-8}{m^3.mol^{-1}}$, soit
\qty{1.49e-3}{cm^3.mol^{-1}}~; $\chi_{\mathrm m}T = 0.445$, $\mu_{\mathrm{eff}} = 1.89\,\mu_B$~: un électron non apparié.
\end{solution}

\begin{solution}{exo:b3:complex-spectra-magnetism:12}
$T_{1/2} = 15\,000/75 = \qty{200}{K}$. À \qty{250}{K}~: $\Delta G = 15\,000 - 250 \times 75 = \qty{-3750}{J/mol}$, $K = \eu^{1.80} = 6.07$,
$x_{\mathrm{HS}} = 0.86$~; $\chi_{\mathrm m}T = 0.86 \times 3.00 = \qty{2.6}{cm^3.K.mol^{-1}}$.
\end{solution}

\begin{solution}{pb:b3:complex-spectra-magnetism:1}
\textbf{1.} [Ar]$3d^3$.
\textbf{2.} $^4$F.
\textbf{3.} $^4$A$_{2g}$ + $^4$T$_{2g}$ + $^4$T$_{1g}$.
\textbf{4.} $^4$A$_{2g}$, le terme de $t_{2g}^3$, avec les trois électrons dans les orbitales les plus basses.
\textbf{5.} $^4$T$_{1g}$(P).
\textbf{6.} $^4$A$_{2g}\to{}^4$T$_{2g}$, $\to{}^4$T$_{1g}$(F), $\to{}^4$T$_{1g}$(P).
\textbf{7.} 561 et \qty{414}{nm}.
\textbf{8.} 597 et \qty{434}{nm}.
\textbf{9.} Le rubis absorbe le jaune-vert et le violet et transmet le rouge
(avec un peu de bleu)~; les bandes de l'émeraude, décalées vers le rouge,
absorbent aussi le rouge orangé, et le vert est transmis.
\textbf{10.} Rubis \qty{17820}{cm^{-1}}, émeraude \qty{16740}{cm^{-1}}.
\textbf{11.} $^4$T$_{2g}$ ($t_{2g}^2e_g$) et $^4$A$_{2g}$ ($t_{2g}^3$) ont la même énergie de répulsion électronique~; leur
différence est l'énergie de promotion monoélectronique $\Delta_{\mathrm o}$.
\textbf{12.} $B = (14\,305 - 547)/15 = \qty{917}{cm^{-1}}$.
\textbf{13.} \qty{614}{cm^{-1}}.
\textbf{14.} \qty{618}{cm^{-1}}.
\textbf{15.} 0.67 pour les deux.
\textbf{16.} 29.0 et 27.1.
\textbf{17.} $\nu_3 = \qty{38500}{cm^{-1}}$ (\qty{260}{nm}) pour le rubis, dans l'ultraviolet, où l'absorption
de transfert de charge la masque.
\textbf{18.} Le site octaédrique du béryl est plus grand~: des liaisons Cr--O
plus longues donnent un éclatement plus petit, qui varie fortement avec la
distance.
\textbf{19.} $10^7/695 = \qty{14390}{cm^{-1}}$, $23.4B$.
\textbf{20.} $^2$E$_g$ dépend surtout de $C$~: le niveau mesuré, plus haut, signifie que $C/B$ est plus grand,
environ 5.3 d'après la même diagonalisation.
\textbf{21.} $^2$E$_g$ et $^4$A$_{2g}$ appartiennent tous deux à $t_{2g}^3$ et ne diffèrent que par le retournement d'un spin~: la liaison,
donc la géométrie, est la même dans les deux, et la transition n'a pas de
progression vibrationnelle (une raie 0--0).
\textbf{22.} Elle est interdite de spin.
\textbf{23.} Son énergie dépend de $B$ et $C$, presque pas de $\Delta_{\mathrm o}$ (une droite horizontale sur le diagramme),
et $B$ est presque le même dans les deux gemmes.
\textbf{24.} Les ions pompés dans les larges bandes se désexcitent et
s'accumulent dans le niveau $^2$E$_g$ de longue durée de vie, si bien que
ce niveau peut contenir plus d'ions que le niveau fondamental~: c'est une
inversion de \omterm{def:b3:partition-functions:boltzmann}{population}.
\textbf{25.} \textbf{$\Delta_{\mathrm o}(\text{rubis}) - \Delta_{\mathrm o}(\text{émeraude}) = \qty{1080}{cm^{-1}}$}~: toute la différence entre le rouge et
le vert.
\end{solution}
